%% Worksheets, self checks and tests, AP Calculus AB Unit 1 (Limits, Continuity, and Differentiability).
%% Written to match the guided notes 1-1 to 1-9 (two-column version). No calculator.
\newcommand{\gr}[2]{\begin{center}\begin{tikzpicture}\begin{axis}[gridaxis,#1]#2\end{axis}\end{tikzpicture}\end{center}}
\newcommand{\pt}[1]{\node[circle,fill=apblue,inner sep=1.5pt] at (axis cs:#1){};}
\newcommand{\op}[1]{\node[circle,draw=apblue,fill=white,thick,inner sep=1.5pt] at (axis cs:#1){};}
\newcommand{\va}[1]{\draw[apred,dashed,thick](axis cs:#1,-100)--(axis cs:#1,100);}
\newcommand{\ha}[1]{\draw[apgold!80!black,dashed,thick](axis cs:-100,#1)--(axis cs:100,#1);}
\newcommand{\work}[1]{\ifkey\else\par\vspace{#1}\fi}
\newcommand{\gfig}[1]{\Needspace{20\baselineskip}\begin{center}\textsf{\small\bfseries\color{apnavy}#1}\end{center}}
\newcommand{\fgraph}{\gr{xmin=-3.5,xmax=4.5,ymin=-0.5,ymax=4.5,width=7.5cm,height=5cm}{\addplot[apblue,very thick,domain=-3:-1]{x+4};\addplot[apblue,very thick,domain=-1:2]{x^2};\addplot[apblue,very thick,domain=2:4]{8-2*x};\op{-1,3}\pt{-1,1}\op{2,4}\pt{2,1}\pt{-3,1}\pt{4,0}\node[font=\small,apblue] at (axis cs:3.6,3.4){$y=f(x)$};}}
\newcommand{\ggraph}{\gr{xmin=-4.5,xmax=5.5,ymin=-4.5,ymax=6.5,width=8cm,height=6.5cm}{\addplot[apblue,very thick,domain=-4:-1]{2};\addplot[apblue,very thick,domain=-1:1]{x+1};\addplot[apblue,very thick,domain=1:2.33,samples=120]{4/(3-x)};\addplot[apblue,very thick,domain=3.34:5,samples=120]{2-2/(x-3)};\va{3}\op{-1,2}\op{-1,0}\pt{-1,3}\op{1,2}\pt{1,4}\pt{-4,2}\pt{5,1}\node[font=\small,apblue,fill=white,inner sep=1pt] at (axis cs:-2.6,3.4){$y=g(x)$};}}
% 1-1: limits from a graph
\newcommand{\gWa}{\gr{xmin=-4.5,xmax=5.5,ymin=-1.5,ymax=4.5,width=8.5cm,height=5.4cm}{\addplot[apblue,very thick] coordinates {(-4,1)(-2,3)};\addplot[apblue,very thick,domain=-2:1,samples=60]{3-(x+2)^2/3};\addplot[apblue,very thick] coordinates {(1,2)(3,4)};\addplot[apblue,very thick] coordinates {(3,1)(5,-1)};\pt{-4,1}\op{-2,3}\pt{-2,1}\pt{1,0}\op{1,2}\pt{3,4}\op{3,1}\pt{5,-1}\node[apblue,font=\small] at (axis cs:4.4,2.6){$y=f(x)$};}}
% 1-4: types of discontinuities
\newcommand{\gWd}{\gr{xmin=-4.5,xmax=5.5,ymin=-5.5,ymax=5.5,width=8.5cm,height=6cm}{\addplot[apblue,very thick] coordinates {(-4,1)(1,3.5)};\op{-2,2}\pt{-2,4}\pt{-4,1}\pt{1,3.5}\addplot[apblue,very thick,domain=1:2.8,samples=60]{1/(3-x)};\op{1,0.5}\addplot[apblue,very thick,domain=3.2:5,samples=60]{-1/(x-3)};\pt{5,-0.5}\va{3}\node[apblue,font=\small] at (axis cs:4.3,3.5){$y=g(x)$};}}
% 1-7: tangent slopes
\newcommand{\gWt}{\gr{xmin=-0.3,xmax=6.8,ymin=-0.5,ymax=5,width=8cm,height=5cm,xtick={0.5,2.057,3.5,4.943,6.2},xticklabels={$a$,$b$,$c$,$d$,$e$},ytick=\empty}{\addplot[apblue,very thick,domain=0:6.5,samples=120]{0.15*(x-1)*(x-3.5)*(x-6)+3};\pt{0.5,1.763}\pt{2.057,3.902}\pt{3.5,3}\pt{4.943,2.098}\pt{6.2,3.421}}}
% 1-7: where a function is not differentiable
\newcommand{\gWf}{\gr{xmin=-4.5,xmax=6.5,ymin=-1.5,ymax=5,width=9.5cm,height=5.6cm,xtick={-4,-3,...,6}}{\addplot[apblue,very thick] coordinates {(-4,2)(-2,-1)(1,2)};\pt{-4,2}\addplot[apblue,very thick,domain=1:4,samples=140]{3.26+sign(x-3)*abs(x-3)^(1/3)};\op{1,2}\pt{4,4.26}\addplot[apblue,very thick] coordinates {(4,1)(6,2)};\op{4,1}\pt{6,2}\node[apblue,font=\small] at (axis cs:-3,4.3){$y=h(x)$};}}

%%DOC 1-1A | Limits to a Point on a Graph
\gfig{Questions 1 to 8 refer to the graph of $f$ on $[-4,5]$.}
\gWa
\begin{exo}{}
What is $\displaystyle\lim_{x\to-2}f(x)$?
\choices{$1$}{$3$}{It does not exist.}{$0$}
\answer{B}{From both sides the graph approaches the open circle at height $3$.}
\end{exo}
\begin{exo}{}
What is $f(-2)$?
\choices{$3$}{$0$}{$1$}{It does not exist.}
\answer{C}{$f(-2)$ is the height of the filled dot at $x=-2$.}
\end{exo}
\begin{exo}{}
What is $\displaystyle\lim_{x\to1^-}f(x)$?
\choices{$0$}{$2$}{It does not exist.}{$1$}
\answer{A}{From the left, the curved piece comes down to the filled dot $(1,0)$.}
\end{exo}
\begin{exo}{}
What is $\displaystyle\lim_{x\to1}f(x)$?
\choices{$0$}{$2$}{$1$}{It does not exist.}
\answer{D}{Left limit $0$, right limit $2$: different, so the limit does not exist.}
\end{exo}
\begin{exo}{}
What is $\displaystyle\lim_{x\to3^+}f(x)$?
\choices{$4$}{$1$}{It does not exist.}{$3$}
\answer{B}{From the right, the graph approaches the open circle at $(3,1)$.}
\end{exo}
\begin{exo}{}
What is $f(3)$?
\choices{$1$}{It does not exist.}{$4$}{$3$}
\answer{C}{The filled dot at $x=3$ is at height $4$.}
\end{exo}
\begin{exo}{}
For which values of $x$ in $(-4,5)$ does $\displaystyle\lim f(x)$ fail to exist?
\choices{$x=1$ and $x=3$}{$x=1$ only}{$x=-2$ and $x=1$}{$x=-2$, $x=1$ and $x=3$}
\answer{A}{At $x=1$ and $x=3$ the one-sided limits differ. At $x=-2$ the limit exists ($3$), even though $f(-2)=1$.}
\end{exo}
\begin{exo}{}
Which statement is true?
\choices{$\displaystyle\lim_{x\to-2}f(x)=f(-2)$}{$\displaystyle\lim_{x\to3^-}f(x)=1$}{$\displaystyle\lim_{x\to4}f(x)=0$}{$f(1)=2$}
\answer{C}{On the last piece, the height at $x=4$ is $0$.}
\end{exo}
\Needspace{14\baselineskip}
\begin{exo}[Free response]{}
\begin{enumerate}[label=(\alph*)]
\item Using the graph of $f$, find $\lim_{x\to-3}f(x)$, $\lim_{x\to3^-}f(x)$ and $\lim_{x\to3}f(x)$.
\item On the grid, sketch a function $g$ with $\lim_{x\to2}g(x)=3$ but $g(2)=1$.
\item Explain the difference between $\lim_{x\to c}f(x)$ and $f(c)$.
\end{enumerate}
\gr{xmin=-1.5,xmax=4.5,ymin=-0.5,ymax=4.5,width=6cm,height=4.5cm}{\ifkey\addplot[apblue,very thick,domain=-1:4]{0.5*x+2};\op{2,3}\pt{2,1}\fi}
\work{4cm}
\begin{solution}
(a) $2$; $4$; does not exist ($4\neq1$). \;(b) Example: a line through $(2,3)$ with an open circle there, and a filled dot at $(2,1)$. \;(c) The limit is the height the graph approaches near $c$; $f(c)$ is the actual height at $c$. They can be different, and either one can fail to exist.
\end{solution}
\end{exo}

%%DOC 1-2A | Limits to a Point of a Function
\begin{exo}{}
What is $\displaystyle\lim_{x\to2}\left(3x^2-x+1\right)$?
\choices{$11$}{$13$}{$9$}{It does not exist.}
\answer{A}{Plug in: $12-2+1=11$.}
\end{exo}
\begin{exo}{}
What is $\displaystyle\lim_{x\to-1}\left(x^3+2x\right)$?
\choices{$3$}{$-3$}{$1$}{$-1$}
\answer{B}{$-1-2=-3$.}
\end{exo}
\begin{exo}{}
What is $\displaystyle\lim_{x\to4}7$?
\choices{$4$}{$0$}{$7$}{$28$}
\answer{C}{A constant function: the limit is the constant.}
\end{exo}
\begin{exo}{}
Let $p(x)=\begin{cases}x^2-3, & x<2\\ 2x-3, & x\ge2.\end{cases}$ What is $\displaystyle\lim_{x\to2}p(x)$?
\choices{It does not exist.}{$1$}{$4$}{$-1$}
\answer{B}{Left: $4-3=1$. Right: $4-3=1$. The limit is $1$.}
\end{exo}
\begin{exo}{}
Let $q(x)=\begin{cases}5-x, & x\le3\\ x^2-6, & x>3.\end{cases}$ What is $\displaystyle\lim_{x\to3}q(x)$?
\choices{$2$}{$3$}{It does not exist.}{$\frac52$}
\answer{C}{Left: $5-3=2$. Right: $9-6=3$. Different: DNE.}
\end{exo}
\begin{exo}{}
For the same $q$, what is $q(3)$?
\choices{$2$}{$3$}{It does not exist.}{$5$}
\answer{A}{$x=3$ belongs to the top piece: $5-3=2$.}
\end{exo}
\begin{exo}{}
What is $\displaystyle\lim_{x\to5}\frac{x^2-25}{x-5}$?
\choices{$0$}{$10$}{$5$}{It does not exist.}
\answer{B}{$\frac00$. $\frac{(x-5)(x+5)}{x-5}=x+5\to10$.}
\end{exo}
\begin{exo}{}
What is $\displaystyle\lim_{x\to-2}\frac{x^2+5x+6}{x^2-4}$?
\choices{$\frac14$}{$0$}{It does not exist.}{$-\frac14$}
\answer{D}{$\frac{(x+2)(x+3)}{(x+2)(x-2)}=\frac{x+3}{x-2}\to\frac{1}{-4}$.}
\end{exo}
\begin{exo}{}
What is $\displaystyle\lim_{x\to0}\frac{2x^2+6x}{x^2-3x}$?
\choices{$-2$}{$2$}{$0$}{It does not exist.}
\answer{A}{$\frac{2x(x+3)}{x(x-3)}=\frac{2(x+3)}{x-3}\to\frac{6}{-3}=-2$.}
\end{exo}
\Needspace{14\baselineskip}
\begin{exo}[Free response]{}
\begin{enumerate}[label=(\alph*)]
\item Find $\displaystyle\lim_{x\to3}\frac{x^2-x-6}{x^2-9}$. Show the three steps.
\item Let $h(x)=\begin{cases}x+4, & x<1\\ x^2+4, & x\ge1.\end{cases}$ Find $\lim_{x\to1^-}h(x)$, $\lim_{x\to1^+}h(x)$, $\lim_{x\to1}h(x)$ and $h(1)$.
\item A classmate says: ``Plugging in gives $\frac00$, so the limit is $0$.'' Explain the mistake.
\end{enumerate}
\work{4cm}
\begin{solution}
(a) $\frac00$. Factor: $\frac{(x-3)(x+2)}{(x-3)(x+3)}$; cancel: $\frac{x+2}{x+3}$; plug in: $\frac56$.\\
(b) $5$, $5$, $5$, $5$.\\
(c) $\frac00$ is an indeterminate form, not a value. We must factor, cancel and plug in again.
\end{solution}
\end{exo}

%%DOC 1-3A | Limits to Infinity and Horizontal Asymptotes
\begin{exo}{}
Which describes the graph of $f(x)=\dfrac{x+1}{x^2-4x-5}$?
\choices{A hole at $x=-1$ and a vertical asymptote at $x=5$}{Vertical asymptotes at $x=-1$ and $x=5$}{A hole at $x=5$ and a vertical asymptote at $x=-1$}{No hole and no vertical asymptote}
\answer{A}{$\frac{x+1}{(x-5)(x+1)}$: $x+1$ cancels (hole), $x-5$ stays (asymptote).}
\end{exo}
\begin{exo}{}
What is the $y$-coordinate of the hole in the graph of $f(x)=\dfrac{x+1}{x^2-4x-5}$?
\choices{$0$}{$-\frac16$}{$\frac16$}{$-\frac14$}
\answer{B}{After cancelling, $\frac{1}{x-5}$ at $x=-1$: $-\frac16$.}
\end{exo}
\begin{exo}{}
For the same $f$, what is $\displaystyle\lim_{x\to5^+}f(x)$?
\choices{$-\infty$}{$0$}{$+\infty$}{$\frac16$}
\answer{C}{$\frac{1}{x-5}$ with $x-5$ small and positive: $+\infty$.}
\end{exo}
\begin{exo}{}
What is $\displaystyle\lim_{x\to\infty}\frac{2x^2-3}{5x^2+x}$?
\choices{$\frac25$}{$0$}{$-3$}{It does not exist.}
\answer{A}{Same degree: $\frac25$.}
\end{exo}
\begin{exo}{}
What is $\displaystyle\lim_{x\to\infty}\frac{8x-1}{2x^3+5}$?
\choices{$4$}{$8$}{It does not exist.}{$0$}
\answer{D}{The bigger degree is on the bottom: $0$.}
\end{exo}
\begin{exo}{}
What is $\displaystyle\lim_{x\to-\infty}\frac{x^3+1}{4x-x^2}$?
\choices{$-1$}{It does not exist (the function grows without bound).}{$0$}{$\frac14$}
\answer{B}{The bigger degree is on top, so there is no horizontal asymptote ($\frac{x^3}{-x^2}=-x\to+\infty$).}
\end{exo}
\begin{exo}{}
Which are all the asymptotes of $y=\dfrac{3x^2+1}{x^2-9}$?
\choices{$x=3$ and $y=3$ only}{$x=-3$, $x=3$ and $y=0$}{$x=-3$, $x=3$ and $y=3$}{$y=3$ only}
\answer{C}{Nothing cancels: $x=\pm3$ are vertical asymptotes. Same degree: HA $y=\frac31=3$.}
\end{exo}
\Needspace{16\baselineskip}
\begin{exo}{}
The graph of $g$ is shown, with its asymptotes. What is $\displaystyle\lim_{x\to-2^-}g(x)$?
\gr{xmin=-6.5,xmax=4.5,ymin=-5,ymax=6,width=7.5cm,height=5.2cm}{\addplot[apblue,very thick,domain=-6.5:-2.4,samples=60]{1+2/(x+2)};\addplot[apblue,very thick,domain=-1.65:4.5,samples=60]{1+2/(x+2)};\va{-2}\ha{1}}
\choices{$1$}{$+\infty$}{$0$}{$-\infty$}
\answer{D}{Just to the left of $x=-2$, the graph goes down without bound.}
\end{exo}
\Needspace{14\baselineskip}
\begin{exo}[Free response]{}
Let $k(x)=\dfrac{2x^2-2}{x^2+2x-3}$.
\begin{enumerate}[label=(\alph*)]
\item Factor $k$. Find the vertical asymptote and the hole (both coordinates).
\item Find $\lim_{x\to\infty}k(x)$ and $\lim_{x\to-\infty}k(x)$. What is the horizontal asymptote?
\item Find $\lim_{x\to-3^+}k(x)$.
\end{enumerate}
\work{4cm}
\begin{solution}
(a) $k=\frac{2(x-1)(x+1)}{(x-1)(x+3)}$. VA $x=-3$; hole at $x=1$, height $\frac{2(2)}{4}=1$: $(1,1)$.\\
(b) Same degree: both limits are $2$; HA $y=2$.\\
(c) $\frac{2(x+1)}{x+3}$: numerator near $-4$, denominator small and positive: $-\infty$.
\end{solution}
\end{exo}

%%DOC 1-4A | Continuity at a Point
\gfig{Questions 1 to 4 refer to the graph of $g$ on $[-4,5]$.}
\gWd
\begin{exo}{}
What type of discontinuity does $g$ have at $x=-2$?
\choices{Removable (hole)}{Jump}{Infinite}{None}
\answer{A}{The limit exists ($2$) but $g(-2)=4$.}
\end{exo}
\begin{exo}{}
What type of discontinuity does $g$ have at $x=1$?
\choices{Removable}{Jump}{Infinite}{None}
\answer{B}{Left limit $3.5$, right limit $0.5$.}
\end{exo}
\begin{exo}{}
What type of discontinuity does $g$ have at $x=3$?
\choices{Removable}{Jump}{Infinite}{None}
\answer{C}{A vertical asymptote: $g\to+\infty$ from the left and $-\infty$ from the right.}
\end{exo}
\begin{exo}{}
At $x=-2$, which part of the continuity test fails?
\choices{The left limit does not equal the right limit.}{$g(-2)$ is not defined.}{The limit is infinite.}{The limit does not equal $g(-2)$.}
\answer{D}{Left $=$ right $=2$, but $g(-2)=4$.}
\end{exo}
\begin{exo}{}
Let $h(x)=\begin{cases}x^2+2, & x\le1\\ 4x-1, & x>1.\end{cases}$ Which is true at $x=1$?
\choices{$h$ is continuous: $3=3=3$.}{Jump discontinuity.}{Removable discontinuity.}{$h(1)$ is undefined.}
\answer{A}{Left $3$, right $3$, $h(1)=3$.}
\end{exo}
\begin{exo}{}
For what value of $k$ is $g(x)=\begin{cases}kx-1, & x<2\\ x^2+k, & x\ge2\end{cases}$ continuous at $x=2$?
\choices{$3$}{$5$}{$-5$}{$1$}
\answer{B}{$2k-1=4+k$, so $k=5$.}
\end{exo}
\begin{exo}{}
For what value of $k$ is $f(x)=\begin{cases}kx^2, & x\le-1\\ 2x+k, & x>-1\end{cases}$ continuous at $x=-1$?
\choices{$1$}{$-1$}{No value of $k$ works.}{Every value of $k$ works.}
\answer{C}{$k=-2+k$ gives $0=-2$: impossible.}
\end{exo}
\begin{exo}{}
Which function is continuous at $x=0$?
\choices{$\dfrac{x^2-x}{x}$}{$\begin{cases}1, & x<0\\ 0, & x\ge0\end{cases}$}{$\dfrac{1}{x^2}$}{$|x|$}
\answer{D}{(A): hole at $0$. (B): jump. (C): infinite.}
\end{exo}
\Needspace{14\baselineskip}
\begin{exo}[Free response]{}
\begin{enumerate}[label=(\alph*)]
\item Let $p(x)=\begin{cases}3x-2, & x\le2\\ x^2, & x>2.\end{cases}$ Prove that $p$ is continuous at $x=2$ with the continuity test.
\item Find $k$ so that $q(x)=\begin{cases}x^2+k, & x<3\\ kx+1, & x\ge3\end{cases}$ is continuous at $x=3$.
\item Find and classify the discontinuities of $r(x)=\dfrac{x^2-1}{x^2-x}$.
\end{enumerate}
\work{5cm}
\begin{solution}
(a) Left $3(2)-2=4$; right $2^2=4$; $p(2)=4$. $4=4=4$: continuous.\\
(b) $9+k=3k+1$, so $k=4$.\\
(c) $r=\frac{(x-1)(x+1)}{x(x-1)}$: removable at $x=1$ (limit $2$), infinite at $x=0$.
\end{solution}
\end{exo}

%%DOC 1-5A | IVT and EVT
\begin{exo}{}
$f$ is continuous on $[-1,3]$, with $f(-1)=5$ and $f(3)=-2$. Which must be true?
\choices{$f(1)=1.5$}{$f(c)=0$ for some $c$ in $(-1,3)$}{The minimum of $f$ is $-2$.}{$f$ is decreasing.}
\answer{B}{IVT: $0$ is between $5$ and $-2$.}
\end{exo}
\begin{exo}{}
$f$ is continuous on $[2,5]$, with $f(2)=3$ and $f(5)=11$. The IVT guarantees that $f(c)$ equals which value for some $c$ in $(2,5)$?
\choices{$12$}{$2$}{$7$}{$0$}
\answer{C}{Only $7$ is between $3$ and $11$.}
\end{exo}
\begin{exo}{}
$h$ is continuous. \;\begin{tabular}{|c|c|c|c|c|c|}\hline$x$&$0$&$1$&$3$&$4$&$6$\\\hline$h(x)$&$-2$&$3$&$1$&$-1$&$2$\\\hline\end{tabular}\; What is the least number of zeros of $h$ on $[0,6]$?
\choices{$1$}{$2$}{$3$}{$4$}
\answer{C}{Sign changes on $[0,1]$, $[3,4]$ and $[4,6]$.}
\end{exo}
\begin{exo}{}
Which function is guaranteed to have an absolute maximum on the given interval?
\choices{$x^3$ on $[0,2]$}{$\frac1x$ on $(0,1]$}{$x^2$ on $(-1,1)$}{$\frac{1}{x-3}$ on $[2,4]$}
\answer{A}{EVT: continuous on a closed interval.}
\end{exo}
\begin{exo}{}
Which interval matches $1\le x<4$?
\choices{$(1,4)$}{$[1,4]$}{$(1,4]$}{$[1,4)$}
\answer{D}{$1$ included (bracket), $4$ excluded (parenthesis).}
\end{exo}
\begin{exo}{}
What does $f(c)\le f(x)$ for every $x$ in $[a,b]$ mean?
\choices{$f$ has an absolute minimum at $x=c$.}{$f$ has an absolute maximum at $x=c$.}{$f(c)=0$.}{$f$ is decreasing.}
\answer{A}{The height at $c$ is below all the others.}
\end{exo}
\begin{exo}{}
$f(x)=\frac1x$ has $f(-1)=-1$ and $f(1)=1$, but $f(x)$ is never $0$. Why doesn't this contradict the IVT?
\choices{$0$ is not between $-1$ and $1$.}{The interval is open.}{$f$ is not continuous on $[-1,1]$.}{The IVT only works for polynomials.}
\answer{C}{$f$ has a vertical asymptote at $x=0$.}
\end{exo}
\Needspace{14\baselineskip}
\begin{exo}{}
The graph of $f$ on the open interval $(0,4)$ is shown. Which is true?
\gr{xmin=-0.5,xmax=4.5,ymin=-0.5,ymax=4.5,width=6cm,height=4.4cm}{\addplot[apblue,very thick,domain=0:4]{0.9*x+0.2};\op{0,0.2}\op{4,3.8}}
\choices{$f$ has a maximum at $x=4$.}{$f$ has a maximum and a minimum on $(0,4)$.}{$f$ has no maximum and no minimum on $(0,4)$.}{The EVT guarantees a maximum.}
\answer{C}{The endpoints are not included, so the highest and lowest heights are never reached. The EVT needs a closed interval.}
\end{exo}
\Needspace{14\baselineskip}
\begin{exo}[Free response]{}
The temperature $T(t)$, in $^\circ$C, is continuous. \;\begin{tabular}{|c|c|c|c|c|}\hline$t$ (h)&$0$&$2$&$5$&$8$\\\hline$T(t)$&$12$&$18$&$15$&$9$\\\hline\end{tabular}
\begin{enumerate}[label=(\alph*)]
\item Must $T(t)=16$ for some $t$ in $(0,2)$? Justify.
\item What is the least number of times $T(t)=16$ on $[0,8]$? Explain.
\item Must $T$ have an absolute maximum on $[0,8]$? What can you say about its value?
\end{enumerate}
\work{4cm}
\begin{solution}
(a) Yes: $T$ is continuous on $[0,2]$ and $12<16<18$ (IVT).\\
(b) At least $2$: once in $(0,2)$ and once in $(2,5)$, since $18>16>15$.\\
(c) Yes, EVT (continuous on the closed interval $[0,8]$). The maximum is at least $18$.
\end{solution}
\end{exo}

%%DOC 1-6A | Average Rate of Change
\begin{exo}{}
What is the average rate of change of $f(x)=x^2-4x$ on $[1,5]$?
\choices{$2$}{$8$}{$-2$}{$4$}
\answer{A}{$f(5)=5$, $f(1)=-3$: $\frac{8}{4}=2$.}
\end{exo}
\begin{exo}{}
What is the average rate of change of $g(x)=x^3$ on $[-1,2]$?
\choices{$9$}{$3$}{$\frac73$}{$1$}
\answer{B}{$\frac{8-(-1)}{3}=3$.}
\end{exo}
\begin{exo}{}
What is the average rate of change of $h(x)=\sqrt x$ on $[4,9]$?
\choices{$5$}{$1$}{$\frac15$}{$\frac{3}{9}$}
\answer{C}{$\frac{3-2}{9-4}=\frac15$.}
\end{exo}
\begin{exo}{}
\begin{tabular}{|c|c|c|c|c|}\hline$t$ (h)&$0$&$3$&$5$&$9$\\\hline$P(t)$ (people)&$120$&$150$&$210$&$300$\\\hline\end{tabular}\; What is the average rate of change of $P$ on $3\le t\le9$?
\choices{$150$ people}{$30$ people per hour}{$20$ people per hour}{$25$ people per hour}
\answer{D}{$\frac{300-150}{9-3}=25$ people per hour.}
\end{exo}
\begin{exo}{}
$C(x)$ is the cost, in FCFA, of $x$ kilograms of rice. What are the units of $\dfrac{C(10)-C(4)}{10-4}$?
\choices{FCFA}{FCFA per kilogram}{kilograms per FCFA}{kilograms}
\answer{B}{Output units per input unit.}
\end{exo}
\Needspace{14\baselineskip}
\begin{exo}{}
The graph of $f$ is shown. What is the average rate of change of $f$ on $[-1,3]$?
\gr{xmin=-1.5,xmax=3.5,ymin=-4.5,ymax=4.5,width=6cm,height=5cm}{\addplot[apblue,very thick,domain=-1.3:3.3,samples=60]{x^2/2-3*x+0.5};\pt{-1,4}\pt{3,-4}}
\choices{$2$}{$-8$}{$-2$}{$-\frac12$}
\answer{C}{$f(-1)=4$, $f(3)=-4$: $\frac{-8}{4}=-2$.}
\end{exo}
\begin{exo}{}
The average rate of change of $f$ on $[0,4]$ is $3$, and $f(0)=5$. What is $f(4)$?
\choices{$17$}{$12$}{$8$}{$7$}
\answer{A}{$\frac{f(4)-5}{4}=3$, so $f(4)=17$.}
\end{exo}
\begin{exo}{}
On a graph, the average rate of change on $[a,b]$ is the slope of
\choices{the tangent line at $a$}{the $x$-axis}{the tangent line at $b$}{the secant line through $(a,f(a))$ and $(b,f(b))$}
\answer{D}{Slope between the two points.}
\end{exo}
\Needspace{14\baselineskip}
\begin{exo}[Free response]{}
A tank is drained. \;\begin{tabular}{|c|c|c|c|c|c|}\hline$t$ (min)&$0$&$10$&$20$&$30$&$40$\\\hline$V(t)$ (L)&$500$&$420$&$360$&$320$&$300$\\\hline\end{tabular}
\begin{enumerate}[label=(\alph*)]
\item Find the average rate of change of $V$ on $[0,40]$. Give units and interpret it.
\item Find the average rate of change of $V$ on $[20,40]$.
\item Find the average rates on $[0,10]$, $[10,20]$, $[20,30]$, $[30,40]$. Is the water leaving faster or slower over time?
\end{enumerate}
\work{4.5cm}
\begin{solution}
(a) $\frac{300-500}{40}=-5$ L/min: on average the tank loses $5$ liters per minute.\\
(b) $\frac{300-360}{20}=-3$ L/min.\\
(c) $-8$, $-6$, $-4$, $-2$ L/min: the rates increase toward $0$, so the water leaves more and more slowly.
\end{solution}
\end{exo}

%%DOC 1-7A | Instantaneous Rates of Change
\gfig{Questions 1 to 3 refer to the graph of $f$.}
\gWt
\begin{exo}{}
At which labeled point is $f'(x)<0$?
\choices{$a$}{$b$}{$c$}{$e$}
\answer{C}{The graph falls at $c$.}
\end{exo}
\begin{exo}{}
At which two labeled points is $f'(x)=0$?
\choices{$b$ and $d$}{$a$ and $e$}{$b$ and $c$}{$c$ and $d$}
\answer{A}{The peak $b$ and the valley $d$ have horizontal tangent lines.}
\end{exo}
\begin{exo}{}
Which statement about $f'(a)$ and $f'(e)$ is true?
\choices{Both are negative.}{$f'(a)>0$ and $f'(e)<0$.}{Both are zero.}{Both are positive.}
\answer{D}{The graph rises at $a$ and at $e$.}
\end{exo}
\begin{exo}{}
$\dfrac{dy}{dx}$ at $x=c$ means
\choices{the average rate of change on $[0,c]$}{the slope of the tangent line at $x=c$}{the value $y(c)$}{the area under the graph}
\answer{B}{Instantaneous rate of change $=$ derivative $=$ slope of the tangent line.}
\end{exo}
\begin{exo}{}
Where is $f(x)=|x+1|$ not differentiable?
\choices{$x=1$}{$x=0$}{$x=-1$}{nowhere}
\answer{C}{Sharp corner at $x=-1$.}
\end{exo}
\gfig{Questions 6 and 7 refer to the graph of $h$ on $[-4,6]$.}
\gWf
\begin{exo}{}
At which $x$ does $h$ have a sharp corner?
\choices{$x=-2$}{$x=1$}{$x=3$}{$x=4$}
\answer{A}{Two straight pieces with different slopes meet at $x=-2$.}
\end{exo}
\begin{exo}{}
Why is $h$ not differentiable at $x=3$?
\choices{It has a hole there.}{It has a jump there.}{It has a sharp corner there.}{It has a vertical tangent line there.}
\answer{D}{The graph is continuous but becomes vertical at $x=3$.}
\end{exo}
\begin{exo}{}
Which statement is true?
\choices{Continuous implies differentiable.}{Differentiable implies continuous.}{A function is differentiable at a jump.}{A function is differentiable at a vertical tangent.}
\answer{B}{If $f$ is differentiable at $c$, it is continuous at $c$; the converse is false.}
\end{exo}
\Needspace{12\baselineskip}
\begin{exo}[Free response]{}
Use the graph of $h$ above. List every $x$ in $[-4,6]$ where $h$ is not differentiable, and give the reason for each.
\work{3.5cm}
\begin{solution}
$x=-4$: endpoint. $x=-2$: sharp corner. $x=1$: hole (discontinuous). $x=3$: vertical tangent. $x=4$: jump (discontinuous). $x=6$: endpoint.
\end{solution}
\end{exo}

%%DOC 1-8A | MVT
\begin{exo}{}
What does the MVT need on $[a,b]$?
\choices{$f$ continuous on $[a,b]$ and differentiable on $(a,b)$}{$f$ increasing}{$f(a)=f(b)$}{$f$ continuous on $(a,b)$ only}
\answer{A}{Both conditions.}
\end{exo}
\begin{exo}{}
$f$ is differentiable, $f(0)=-2$ and $f(4)=10$. The MVT guarantees a $c$ in $(0,4)$ with $f'(c)=$
\choices{$12$}{$3$}{$2$}{$\frac13$}
\answer{B}{$\frac{10-(-2)}{4}=3$.}
\end{exo}
\begin{exo}{}
For which function does the MVT apply on $[0,2]$?
\choices{$|x-1|$}{$\frac{1}{x-1}$}{$x^2-3x$}{$\begin{cases}x, & x<1\\ x+1, & x\ge1\end{cases}$}
\answer{C}{A polynomial. The others have a corner, an asymptote or a jump at $x=1$.}
\end{exo}
\begin{exo}{}
$h$ is differentiable. \;\begin{tabular}{|c|c|c|c|}\hline$x$&$1$&$4$&$6$\\\hline$h(x)$&$2$&$11$&$5$\\\hline\end{tabular}\; The MVT guarantees a $c$ in $(4,6)$ with $h'(c)=$
\choices{$3$}{$6$}{$-6$}{$-3$}
\answer{D}{$\frac{5-11}{6-4}=-3$.}
\end{exo}
\begin{exo}{}
$f$ is differentiable and $f(1)=f(5)=3$. What does the MVT guarantee?
\choices{A $c$ in $(1,5)$ with $f'(c)=0$ (a horizontal tangent)}{$f(3)=3$}{$f$ is constant.}{A $c$ in $(1,5)$ with $f(c)=0$}
\answer{A}{Average rate $\frac{3-3}{4}=0$.}
\end{exo}
\Needspace{14\baselineskip}
\begin{exo}{}
The graph of $g(x)=|x-2|+1$ on $[0,4]$ is shown. Does the MVT guarantee a $c$ in $(0,4)$ with $g'(c)=\frac{g(4)-g(0)}{4-0}$?
\gr{xmin=-0.5,xmax=4.5,ymin=-0.5,ymax=3.5,width=6cm,height=4.4cm}{\addplot[apblue,very thick] coordinates {(0,3)(2,1)(4,3)};\pt{0,3}\pt{4,3}}
\choices{Yes, $c=2$.}{No: $g$ is not differentiable at $x=2$.}{Yes, because $g$ is continuous.}{No, because $g(0)=g(4)$.}
\answer{B}{The corner at $x=2$ breaks the conditions; the slopes are only $-1$ and $1$, never $0$.}
\end{exo}
\begin{exo}{}
In words, the MVT says that somewhere the slope of a \underline{\hspace{1.5cm}} equals the slope of the \underline{\hspace{1.5cm}}.
\choices{secant line; tangent line}{tangent line; $x$-axis}{tangent line; secant line}{graph; table}
\answer{C}{Instantaneous rate (tangent) $=$ average rate (secant).}
\end{exo}
\Needspace{14\baselineskip}
\begin{exo}[Free response]{}
The water level $W(t)$, in meters, is differentiable. \;\begin{tabular}{|c|c|c|c|c|}\hline$t$ (h)&$0$&$4$&$6$&$10$\\\hline$W(t)$&$2$&$5$&$9$&$7$\\\hline\end{tabular}
\begin{enumerate}[label=(\alph*)]
\item Explain why there must be a time $t$ in $(4,6)$ with $W'(t)=2$.
\item Must there be a time in $(0,10)$ when the level rises at $0.5$ m/h? Justify.
\item Does the MVT apply to $g(x)=|x-3|+1$ on $[1,5]$? Explain.
\end{enumerate}
\work{4.5cm}
\begin{solution}
(a) $W$ is differentiable (so continuous) on $[4,6]$; $\frac{9-5}{2}=2$: MVT.\\
(b) Yes: $\frac{7-2}{10}=0.5$, MVT on $[0,10]$.\\
(c) No: $g$ has a sharp corner at $x=3$, so it is not differentiable on $(1,5)$.
\end{solution}
\end{exo}

%%DOC 1-9A | IVT, EVT, and MVT
\begin{exo}{}
A runner's distance $D(t)$ is differentiable. ``Was the runner ever going exactly as fast as her average speed?'' Which theorem answers this?
\choices{IVT}{EVT}{MVT}{None}
\answer{C}{A speed (slope) equal to an average rate: MVT.}
\end{exo}
\begin{exo}{}
The height $h(t)$ of a kite is continuous on $[0,30]$. ``Was the kite ever exactly $25$ m high?'' given $h(0)=2$ and $h(30)=40$. Which theorem answers this?
\choices{IVT}{EVT}{MVT}{None}
\answer{A}{A height between $h(0)$ and $h(30)$: IVT.}
\end{exo}
\begin{exo}{}
``Was there a highest point of the kite on $[0,30]$?'' Which theorem answers this?
\choices{IVT}{MVT}{None}{EVT}
\answer{D}{Continuous on a closed interval: an absolute maximum exists (EVT).}
\end{exo}
\begin{exo}{}
Which theorem needs a \textbf{closed} interval and gives an absolute maximum and minimum?
\choices{IVT}{EVT}{MVT}{All three}
\answer{B}{The EVT.}
\end{exo}
\gfig{Questions 5 to 7 refer to the table. $f$ is differentiable.}
\begin{center}\begin{tabular}{|c|c|c|c|c|}\hline$x$&$0$&$3$&$5$&$8$\\\hline$f(x)$&$2$&$-4$&$6$&$6$\\\hline\end{tabular}\end{center}
\begin{exo}{}
Which statement is guaranteed?
\choices{$f(c)=0$ for some $c$ in $(5,8)$}{$f(c)=0$ for some $c$ in $(0,3)$}{$f'(c)=0$ for some $c$ in $(0,3)$}{$f'(c)=-2$ for some $c$ in $(3,5)$}
\answer{B}{IVT: $0$ is between $2$ and $-4$.}
\end{exo}
\begin{exo}{}
The MVT guarantees a $c$ in $(3,5)$ with $f'(c)=$
\choices{$1$}{$-5$}{$5$}{$10$}
\answer{C}{$\frac{6-(-4)}{2}=5$.}
\end{exo}
\begin{exo}{}
Must there be a $c$ in $(5,8)$ with $f'(c)=0$?
\choices{Yes, by the MVT.}{Yes, by the IVT.}{No, $f$ is constant there.}{No, the MVT needs $f(a)\neq f(b)$.}
\answer{A}{$\frac{6-6}{3}=0$ and $f$ is differentiable.}
\end{exo}
\begin{exo}{}
$h$ is continuous on $[0,4]$ with a sharp corner at $x=2$. Which theorems apply on $[0,4]$?
\choices{All three}{MVT only}{IVT and EVT only}{None}
\answer{C}{The MVT also needs differentiability.}
\end{exo}
\Needspace{14\baselineskip}
\begin{exo}[Free response]{}
$g$ is differentiable. \;\begin{tabular}{|c|c|c|c|c|}\hline$x$&$1$&$2$&$4$&$7$\\\hline$g(x)$&$6$&$1$&$5$&$11$\\\hline\end{tabular}
\begin{enumerate}[label=(\alph*)]
\item Must $g(c)=3$ for some $c$ in $(1,2)$? Justify.
\item Must $g'(c)=2$ for some $c$ in $(2,4)$? Justify.
\item Must $g$ have an absolute minimum on $[1,7]$? Justify.
\item $k(x)=\frac{1}{x-3}$ has $k(2)=-1$ and $k(4)=1$. Does the IVT guarantee $k(c)=0$ in $(2,4)$? Explain.
\end{enumerate}
\work{5cm}
\begin{solution}
(a) Yes, IVT: $g$ is continuous (differentiable) and $1<3<6$.\\
(b) Yes, MVT: $\frac{5-1}{4-2}=2$.\\
(c) Yes, EVT: continuous on the closed interval $[1,7]$; the minimum is at most $1$.\\
(d) No: $k$ is not continuous at $x=3$ (vertical asymptote), so the IVT does not apply.
\end{solution}
\end{exo}

%%DOC SC1.1 | Self Check 1.1 \textbullet\ Lessons 1-1 and 1-2
\gfig{Questions 1 and 2 refer to the graph of $f$.}
\gWa
\begin{exo}{}
Using the graph of $f$, $\displaystyle\lim_{x\to-3}f(x)=$
\choices{$1$}{$2$}{$3$}{It does not exist.}
\answer{B}{On the first piece, the height at $x=-3$ is $2$.}
\end{exo}
\begin{exo}{}
Using the graph of $f$, $\displaystyle\lim_{x\to3^-}f(x)=$
\choices{$1$}{$3$}{It does not exist.}{$4$}
\answer{D}{From the left, the graph rises to the filled dot $(3,4)$.}
\end{exo}
\begin{exo}{}
$\displaystyle\lim_{x\to1}\frac{x^2+2x-3}{x-1}=$
\choices{$4$}{$0$}{$2$}{It does not exist.}
\answer{A}{$\frac{(x-1)(x+3)}{x-1}=x+3\to4$.}
\end{exo}
\begin{exo}{}
Let $j(x)=\begin{cases}2x+3, & x<-1\\ x^2, & x\ge-1.\end{cases}$ Then $\displaystyle\lim_{x\to-1}j(x)=$
\choices{$-1$}{It does not exist.}{$1$}{$3$}
\answer{C}{Left: $-2+3=1$. Right: $(-1)^2=1$.}
\end{exo}
\begin{exo}[Free response]{}
Find $\displaystyle\lim_{x\to0}\frac{x^2+5x}{x}$, $\displaystyle\lim_{x\to-3}\frac{x^2-9}{x+3}$, and $\displaystyle\lim_{x\to2}j(x)$ for $j(x)=\begin{cases}x^2, & x<2\\ 6-x, & x\ge2.\end{cases}$
\work{3cm}
\begin{solution}
$x+5\to5$; \;$x-3\to-6$; \;left $4$, right $4$: the limit is $4$.
\end{solution}
\end{exo}

%%DOC SC1.2 | Self Check 1.2 \textbullet\ Lessons 1-3 and 1-4
\begin{exo}{}
$\displaystyle\lim_{x\to\infty}\frac{3x^2+1}{x^2-5}=$
\choices{$3$}{$-\frac15$}{$0$}{$\infty$}
\answer{A}{}
\end{exo}
\begin{exo}{}
$\displaystyle\lim_{x\to2}\frac{3}{(x-2)^2}=$
\choices{$\infty$}{$-\infty$}{$3$}{It does not exist and is not infinite.}
\answer{A}{}
\end{exo}
\begin{exo}{}
For what $a$ is $f(x)=\begin{cases}ax+2, & x<3\\ x^2-a, & x\ge3\end{cases}$ continuous?
\choices{$\frac{11}4$}{$\frac74$}{$\frac72$}{$7$}
\answer{B}{$3a+2=9-a$.}
\end{exo}
\begin{exo}{}
Which condition for continuity fails for $f(x)=\frac{x^2-9}{x-3}$ at $x=3$?
\choices{$f(3)$ is defined}{the limit exists}{the limit equals $f(3)$}{none}
\answer{A}{$f(3)$ is undefined.}
\end{exo}
\begin{exo}[Free response]{}
Let $f(x)=\dfrac{2x^2-8}{x^2-x-2}$. Find the vertical asymptote, the hole and the horizontal asymptote. Then classify the discontinuities of $g(x)=\frac{x+2}{x^2-4}$.
\work{3.5cm}
\begin{solution}
$f=\frac{2(x-2)(x+2)}{(x-2)(x+1)}$: VA $x=-1$; hole at $x=2$, height $\frac{2(4)}{3}=\frac83$; HA $y=2$. \;$g$: removable at $x=-2$ (limit $-\frac14$), infinite at $x=2$.
\end{solution}
\end{exo}

%%DOC SC1.1-1.2 | Self Check 1.1 and 1.2 \textbullet\ Lessons 1-1 to 1-4
\begin{exo}{}
$\displaystyle\lim_{x\to-3}\frac{x^2+x-6}{x^2-9}=$
\choices{$\frac56$}{$0$}{$-\frac56$}{It does not exist.}
\answer{A}{$\frac{(x+3)(x-2)}{(x+3)(x-3)}=\frac{x-2}{x-3}\to\frac{-5}{-6}$.}
\end{exo}
\begin{exo}{}
$\displaystyle\lim_{x\to\infty}\frac{2x-x^3}{4x^3+1}=$
\choices{$2$}{$-\frac14$}{$0$}{It does not exist.}
\answer{B}{Same degree: $\frac{-1}{4}$.}
\end{exo}
\begin{exo}{}
Let $j(x)=\begin{cases}x+3, & x<1\\ 2x, & x\ge1.\end{cases}$ Then $\displaystyle\lim_{x\to1}j(x)=$
\choices{$2$}{$4$}{$3$}{It does not exist.}
\answer{D}{Left $4$, right $2$.}
\end{exo}
\gfig{Question 4 and the free response refer to the graph of $g$.}
\gWd
\begin{exo}{}
Using the graph of $g$, at which $x$ is the discontinuity removable?
\choices{$x=-2$}{$x=1$}{$x=3$}{$x=-4$}
\answer{A}{The limit at $x=-2$ is $2$, but $g(-2)=4$.}
\end{exo}
\begin{exo}[Free response]{}
Using the graph of $g$, classify the discontinuities at $x=-2$, $1$ and $3$, with limit statements.
\work{3cm}
\begin{solution}
$-2$: removable ($\lim=2\neq4=g(-2)$). $1$: jump (left $3.5$, right $0.5$). $3$: infinite ($+\infty$ from the left, $-\infty$ from the right).
\end{solution}
\end{exo}

%%DOC SC1.3 | Self Check 1.3 \textbullet\ Lessons 1-5 and 1-6
\begin{exo}{}
$f$ is continuous, $f(0)=3$, $f(4)=-1$. Which must be true?
\choices{$f(2)=1$}{$f(c)=0$ for some $c$ in $(0,4)$}{exactly one zero}{$f(c)=4$}
\answer{B}{}
\end{exo}
\begin{exo}{}
Why does $\frac1x$ on $[-1,1]$ never equal $0$, although $f(-1)<0<f(1)$?
\choices{$0$ is not between them}{IVT is only for polynomials}{not continuous on $[-1,1]$}{it is decreasing}
\answer{C}{}
\end{exo}
\begin{exo}{}
What is the average rate of change of $h(x)=\frac{12}{x}$ on $[2,6]$?
\choices{$-1$}{$-4$}{$1$}{$-\frac12$}
\answer{A}{}
\end{exo}
\begin{exo}[Free response]{}
Find the absolute maximum and minimum of $f(x)=x^2-4x+5$ on $[0,3]$, and explain why they must exist. Then find the average rate of change of $f$ on $[0,3]$.
\work{3.5cm}
\begin{solution}
Continuous on a closed interval (EVT). Max $5$ at $x=0$, min $1$ at $x=2$. Average rate $\frac{2-5}{3}=-1$.
\end{solution}
\end{exo}

%%DOC SC1.4 | Self Check 1.4 \textbullet\ Lessons 1-7 to 1-9
\begin{exo}{}
$f$ is differentiable, with $f(0)=1$ and $f(3)=10$. The MVT guarantees a $c$ in $(0,3)$ with $f'(c)=$
\choices{$9$}{$3$}{$\frac{10}{3}$}{$\frac13$}
\answer{B}{$\frac{10-1}{3-0}=3$.}
\end{exo}
\begin{exo}{}
Where is $f(x)=\sqrt[3]{x}$ continuous but not differentiable?
\choices{$x=1$}{$x=0$}{$x=-1$}{nowhere}
\answer{B}{Vertical tangent.}
\end{exo}
\begin{exo}{}
$f$ is continuous, $f(1)=2$ and $f(4)=8$. Which theorem guarantees $f(c)=5$ for some $c$ in $(1,4)$?
\choices{MVT}{EVT}{IVT}{None of them}
\answer{C}{A height between $f(1)$ and $f(4)$: the IVT.}
\end{exo}
\begin{exo}[Free response]{}
A differentiable function has $f(1)=2$, $f(3)=8$, $f(6)=-1$. Which theorem guarantees each, and why? (a) $f(c)=5$ for some $c$ in $(1,3)$. (b) $f'(c)=3$ for some $c$ in $(1,3)$. (c) $f'(c)=-3$ for some $c$ in $(3,6)$.
\work{3.5cm}
\begin{solution}
(a) IVT ($2<5<8$, continuous). (b) MVT ($\frac{8-2}{2}=3$). (c) MVT ($\frac{-1-8}{3}=-3$).
\end{solution}
\end{exo}

%%DOC SC1.3-1.4 | Self Check 1.3 and 1.4 \textbullet\ Lessons 1-5 to 1-9
\begin{exo}{}
Which function must have an absolute minimum on the interval?
\choices{$\frac1x$ on $[1,3]$}{$\frac1x$ on $(0,3]$}{$x$ on $(0,1)$}{$\frac{1}{x-2}$ on $[1,3]$}
\answer{A}{}
\end{exo}
\begin{exo}{}
What is the average rate of change of $g(x)=x^2-2x$ on $[1,4]$?
\choices{$2$}{$3$}{$9$}{$-3$}
\answer{B}{$\frac{8-(-1)}{4-1}=3$.}
\end{exo}
\begin{exo}{}
$h$ is continuous on $[0,5]$ and its graph has a sharp corner at $x=2$. Which theorems are guaranteed to apply on $[0,5]$?
\choices{MVT only}{IVT, EVT and MVT}{IVT and EVT only}{None of them}
\answer{C}{The MVT also needs differentiability, which fails at the corner.}
\end{exo}
\begin{exo}[Free response]{}
$f$ is differentiable. \;\begin{tabular}{|c|c|c|c|c|}\hline$x$&$0$&$2$&$5$&$7$\\\hline$f(x)$&$3$&$-1$&$8$&$8$\\\hline\end{tabular}\; Justify each answer with a theorem. (a) Must $f(c)=0$ for some $c$ in $(0,2)$? (b) Must $f'(c)=3$ for some $c$ in $(2,5)$? (c) Must $f'(c)=0$ for some $c$ in $(5,7)$? (d) Must $f$ have an absolute maximum on $[0,7]$?
\work{3.5cm}
\begin{solution}
(a) Yes, IVT: continuous (because differentiable) and $0$ is between $3$ and $-1$. (b) Yes, MVT: $\frac{8-(-1)}{3}=3$. (c) Yes, MVT: $\frac{8-8}{2}=0$. (d) Yes, EVT: continuous on the closed interval $[0,7]$.
\end{solution}
\end{exo}

%%DOC PRETEST1 | Unit 1 Pretest \textbullet\ Limits, Continuity, and Differentiability
\partheading{Part A \textbullet\ Multiple choice}
\begin{exo}{}
$\displaystyle\lim_{x\to3}\frac{x^2-9}{x-3}$ is
\choices{$0$}{$3$}{$6$}{nonexistent}
\answer{C}{}
\end{exo}
\begin{exo}{}
$\displaystyle\lim_{x\to\infty}\frac{5x^2-x}{2x^2+3}$ is
\choices{$0$}{$\frac52$}{$5$}{$\infty$}
\answer{B}{}
\end{exo}
\begin{exo}{}
$\displaystyle\lim_{x\to1^-}\frac{2}{x-1}$ is
\choices{$\infty$}{$2$}{$0$}{$-\infty$}
\answer{D}{}
\end{exo}
\begin{exo}{}
$\displaystyle\lim_{x\to2}\frac{x^2+x-6}{x-2}$ is
\choices{$0$}{$1$}{$5$}{nonexistent}
\answer{C}{$\frac{(x-2)(x+3)}{x-2}=x+3\to5$.}
\end{exo}
\fgraph
\begin{exo}{}
Using the graph of $f$, which limit does not exist?
\choices{$\lim_{x\to0}f(x)$}{$\lim_{x\to-1}f(x)$}{$\lim_{x\to2}f(x)$}{$\lim_{x\to3}f(x)$}
\answer{B}{}
\end{exo}
\begin{exo}{}
Using the graph of $f$, at which $x$ is the discontinuity removable?
\choices{$x=-1$}{$x=0$}{$x=2$}{$x=3$}
\answer{C}{}
\end{exo}
\begin{exo}{}
For which $k$ is $f(x)=\begin{cases}kx^2, & x\le1\\ x+3, & x>1\end{cases}$ continuous?
\choices{$4$}{$3$}{$1$}{$2$}
\answer{A}{}
\end{exo}
\begin{exo}{}
$f$ is continuous on $[2,6]$, $f(2)=5$, $f(6)=-3$. Which must be true?
\choices{$f(4)=1$}{$f(c)=0$ for some $c$}{$f$ has no maximum}{$f'(c)=-2$}
\answer{B}{IVT. (D) needs differentiability.}
\end{exo}
\begin{exo}{}
$f$ is differentiable, $f(1)=2$ and $f(5)=10$. Which must be true?
\choices{$f'(c)=2$ for some $c$ in $(1,5)$}{$f'(c)=8$ for some $c$ in $(1,5)$}{$f(3)=6$}{$f'(c)=0$ for some $c$ in $(1,5)$}
\answer{A}{MVT: $\frac{10-2}{5-1}=2$.}
\end{exo}
\begin{exo}{}
$f(x)=|x+1|$ is not differentiable at
\choices{$x=1$}{$x=0$}{$x=-1$}{nowhere}
\answer{C}{}
\end{exo}
\partheading{Part B \textbullet\ Free response}
\begin{exo}[Free response]{}
Let $g(x)=\begin{cases}\dfrac{x^2+x-6}{x-2}, & x<2\\[4pt] cx+1, & x\ge2.\end{cases}$
\begin{enumerate}[label=(\alph*)]
\item Find $\lim_{x\to2^-}g(x)$.
\item Find $c$ so that $g$ is continuous at $x=2$; justify with the definition.
\item With this $c$, is $g$ differentiable at $x=2$? (The left piece simplifies to $x+3$.)
\end{enumerate}
\work{3.5cm}
\begin{solution}
(a) $x+3\to5$. (b) $2c+1=5$, $c=2$; $g(2)=5=\lim$. (c) Left slope $1$, right slope $2$: not differentiable (corner).
\end{solution}
\end{exo}
\begin{exo}[Free response]{}
A differentiable function $W$ (temperature, $^\circ$C) has $W(0)=20$, $W(3)=24$, $W(5)=31$, $W(8)=27$.
\begin{enumerate}[label=(\alph*)]
\item Must $W(t)=30$ for some $t$ in $(3,5)$? Justify.
\item Must $W'(t)=3.5$ for some $t$ in $(3,5)$? Justify.
\item Must $W$ have an absolute maximum on $[0,8]$? Justify.
\end{enumerate}
\work{3.5cm}
\begin{solution}
(a) Yes, IVT ($24<30<31$; differentiable implies continuous). (b) Yes, MVT ($\frac{31-24}{2}=3.5$). (c) Yes, EVT: $W$ is continuous on the closed interval $[0,8]$; the maximum is at least $31$.
\end{solution}
\end{exo}

%%DOC FIXUP1 | Unit 1 Fix-Up \textbullet\ Limits, Continuity, and Differentiability
\partheading{Part A \textbullet\ Multiple choice}
\begin{exo}{}
$\displaystyle\lim_{x\to-4}\frac{x^2-16}{x+4}$ is
\choices{$0$}{$-8$}{$8$}{nonexistent}
\answer{B}{}
\end{exo}
\begin{exo}{}
$\displaystyle\lim_{x\to\infty}\frac{2x-7}{x^2+1}$ is
\choices{$0$}{$2$}{$-7$}{$\infty$}
\answer{A}{}
\end{exo}
\begin{exo}{}
$\displaystyle\lim_{x\to3^+}\frac{-1}{x-3}$ is
\choices{$\infty$}{$-1$}{$-\infty$}{$0$}
\answer{C}{}
\end{exo}
\begin{exo}{}
$\displaystyle\lim_{x\to-2}\frac{x^2+5x+6}{x+2}$ is
\choices{$0$}{$1$}{$-1$}{nonexistent}
\answer{B}{$\frac{(x+2)(x+3)}{x+2}=x+3\to1$.}
\end{exo}
\ggraph
\begin{exo}{}
Using the graph of $g$, which limit does not exist?
\choices{$\lim_{x\to1}g(x)$}{$\lim_{x\to4}g(x)$}{$\lim_{x\to-1}g(x)$}{$\lim_{x\to-2}g(x)$}
\answer{C}{}
\end{exo}
\begin{exo}{}
Using the graph of $g$, what type of discontinuity is at $x=3$?
\choices{removable}{jump}{infinite}{none}
\answer{C}{}
\end{exo}
\begin{exo}{}
For which $k$ is $f(x)=\begin{cases}2x+k, & x<1\\ kx^2+3, & x\ge1\end{cases}$ continuous?
\choices{$1$}{$-1$}{It is continuous for every $k$.}{No value of $k$ works.}
\answer{D}{$2+k=k+3$ is impossible.}
\end{exo}
\begin{exo}{}
Which function must have an absolute maximum on $[0,2]$?
\choices{$\frac{1}{x-1}$}{$\frac{1}{x+1}$}{$\frac{1}{x-2}$}{$\frac{1}{x}$}
\answer{B}{Continuous on the closed interval $[0,2]$ (EVT). The others have a vertical asymptote in $[0,2]$.}
\end{exo}
\begin{exo}{}
Where is $f(x)=\sqrt[3]{x-1}$ not differentiable?
\choices{$x=1$}{$x=0$}{$x=-1$}{nowhere}
\answer{A}{Vertical tangent line at $x=1$.}
\end{exo}
\begin{exo}{}
$f$ is differentiable with $f(2)=4$ and $f(6)=36$. The MVT guarantees a $c$ in $(2,6)$ with $f'(c)=$
\choices{$3$}{$5$}{$8$}{$32$}
\answer{C}{$\frac{36-4}{6-2}=8$.}
\end{exo}
\partheading{Part B \textbullet\ Free response}
\begin{exo}[Free response]{}
Let $h(x)=\dfrac{x^2-x-2}{x^2-1}$.
\begin{enumerate}[label=(\alph*)]
\item Find and classify each discontinuity, with limits.
\item Find $\lim_{x\to\pm\infty}h(x)$.
\item What value of $h$ at the removable discontinuity makes $h$ continuous there?
\end{enumerate}
\work{3.5cm}
\begin{solution}
$h=\frac{(x-2)(x+1)}{(x-1)(x+1)}$. (a) $x=-1$: removable ($\lim=\frac{-3}{-2}=\frac32$). $x=1$: infinite ($\lim_{x\to1^-}=+\infty$, $\lim_{x\to1^+}=-\infty$). (b) $1$. (c) $\frac32$.
\end{solution}
\end{exo}
\begin{exo}[Free response]{}
Let $f(x)=x^3-4x+1$.
\begin{enumerate}[label=(\alph*)]
\item Use the IVT to show that $f$ has a zero in $(0,1)$.
\item Find the average rate of change of $f$ on $[0,2]$.
\item Explain why there is a $c$ in $(0,2)$ where the tangent line to the graph of $f$ is horizontal.
\end{enumerate}
\work{3.5cm}
\begin{solution}
(a) $f$ is a polynomial, so it is continuous on $[0,1]$; $f(0)=1>0>-2=f(1)$, so the IVT gives a zero in $(0,1)$.\\
(b) $\frac{f(2)-f(0)}{2-0}=\frac{1-1}{2}=0$.\\
(c) $f$ is differentiable (a polynomial), so the MVT gives $c$ in $(0,2)$ with $f'(c)=0$: a horizontal tangent.
\end{solution}
\end{exo}
